\documentclass{article}

\usepackage[a4paper, top    = 2.0cm, 
                     bottom = 2.0cm, 
                     left   = 2.5cm, 
                     right  = 2.0cm]{geometry}

\usepackage[brazil]{babel}
\usepackage[T1]{fontenc}
    \DeclareFontShape{T1}{cmr}{m}{scit}{<->ssub * cmr/m/scsl}{}
    \DeclareFontShape{T1}{cmr}{bx}{scit}{<->ssub * cmr/bx/scsl}{}

\usepackage{indentfirst}
\usepackage{float}
\usepackage{graphicx}
    \graphicspath{{figures/}}
    
\usepackage{amsmath}

\usepackage{url}
\usepackage{hyperref}

\usepackage{m_template}

\title{PSI3441 - Arquitetura de Sistemas Embarcados}
\subtitle{Atividade 1}
\author{Gabriel Ivo Bozi & 14571845}

\begin{document}

\maketitle

\section{Por que os \textit{leds} são \textit{Active Low} (acendem quando se coloca 0 na saída)?}

    Na eletrônica dos microcontroladores, os pinos de saída geralmente possuem uma capacidade consideravelmente maior de drenar corrente para o terra (\textit{current sink}) do que de fornecer corrente a partir da fonte de alimentação (\textit{current source}). 
    
    Portanto ligar os \textit{leds} na configuração \textit{Active Low} reduz o estresse elétrico no \textit{chip} e previne danos físicos às portas do microcontrolador. Além disso, essa configuração garante estabilidade, evitando que o \textit{led} pisque ou acenda acidentalmente devido a ruídos no barramento durante o processo de inicialização do sistema, momento em que os pinos frequentemente assumem um estado de alta impedância.

\section{Quais funções você usou para acender e apagar os \textit{leds}?}

    Para configuração e verificação de como os \textit{leds} devem se comportar foram utilizadas as seguintes funções:
    
    \begin{description}
        \item[\texttt{gpio\_pin\_configure\_dt()}] Configura os pinos dos \textit{leds} como saída e define o estado padrão.
        
        \item[\texttt{gpio\_is\_ready\_dt()}] Utilizada como etapa de validação de segurança para confirmar se o periférico mapeado está pronto para receber comandos do sistema.
    \end{description}
    
    Para acionamento dos \textit{leds} foram utilizadas as seguintes funções:
    
    \begin{description}
        \item[\texttt{gpio\_pin\_toggle()}] Inverte o nível lógico no pino.
    
        \item[\texttt{gpio\_pin\_set\_dt()}] Utilizada para acender e apagar os \textit{leds} de fato, alterando o estado lógico do pino (o valor passado como parâmetro representa o estado lógico desejado para o pino, sendo interpretado conforme a configuração do hardware, por exemplo, \textit{Active Low} ou \textit{Active High}).. 
    \end{description}

\section{Explique o que é o \textit{Device Tree}}
    O \textit{DeviceTree} é uma estrutura de dados que descreve o \textit{hardware} da plataforma de forma declarativa. Em vez de utilizar endereços ou identificadores fixos no código, ele permite referenciar periféricos por meio de aliases e propriedades, com isso um pino com uma função especifica, por exemplo o \textit{led}, pode funcionar em diferentes plataformas com arranjos dos pinos diferentes. , e é capaz de descrever propriedades dos pinos, por exemplo a polaridade do \textit{led} já que nessa placa ele é ativo em nível lógico baixo. Além de ser capaz de fazer verificações de segurança e portabilidade em tempo de compilação.

\section{Explique as abstrações feitas pelo Sistema Operacional}
    Sistemas operacionais de tempo real como o \textit{Zephyr} constroem camadas de abstração de hardware para isolar o desenvolvedor da manipulação direta de registradores e mapeamento da memória por meio do \textit{linker}.

\appendix
\renewcommand{\thesection}{\arabic{section}}
\addcontentsline{toc}{section}{Anexos}  
    
\newpage
\part*{\centering Anexos}
    \section{Código main.c utilizado}
        \begin{small}
        \begin{verbatim}
#include <zephyr/kernel.h>
#include <zephyr/device.h>
#include <zephyr/drivers/gpio.h>

#define SLEEP_TIME_MS 500

// Define os LEDs usando Device Tree Alias
#define LED0_NODE DT_ALIAS(led0) // green
#define LED1_NODE DT_ALIAS(led1) // blue
#define LED2_NODE DT_ALIAS(led2) // red

#if DT_NODE_HAS_STATUS(LED0_NODE, okay) || DT_NODE_HAS_STATUS(LED1_NODE, okay) || 
            DT_NODE_HAS_STATUS(LED2_NODE, okay)
    static const struct gpio_dt_spec led0 = GPIO_DT_SPEC_GET(LED0_NODE, gpios);
    static const struct gpio_dt_spec led1 = GPIO_DT_SPEC_GET(LED1_NODE, gpios);
    static const struct gpio_dt_spec led2 = GPIO_DT_SPEC_GET(LED2_NODE, gpios);    
#else
    #error "Unsupported board: leds device tree alias is not defined"
#endif

// Definição dos estados possíveis para a máquina de estados
//  red
//  yellow = green + red
//  green
typedef enum 
{
    STATE_RED,
    STATE_YELLOW,
    STATE_GREEN
} led_state_t;

void main()
{
    // Verifica se os devices estão prontos
    if (!gpio_is_ready_dt(&led0) || !gpio_is_ready_dt(&led1) || !gpio_is_ready_dt(&led2)) 
    {
        printk("Error: One or more leds devices are not ready\n");
        return;
    }

    // Configura os 3 pinos como saída e inicia desligados
    gpio_pin_configure_dt(&led0, GPIO_OUTPUT_INACTIVE);
    gpio_pin_configure_dt(&led1, GPIO_OUTPUT_INACTIVE);
    gpio_pin_configure_dt(&led2, GPIO_OUTPUT_INACTIVE);

    printk("leds ready to blink using a state machine\n");

    // Inicializa a máquina de estados no verde
    led_state_t current_state = STATE_GREEN;

    while (1) 
    {
        // Desliga todos os leds
        gpio_pin_set_dt(&led0, 0);
        gpio_pin_set_dt(&led1, 0);
        gpio_pin_set_dt(&led2, 0);

        switch (current_state) 
        {
            //
            case STATE_GREEN: 
                gpio_pin_set_dt(&led0, 1);
                current_state = STATE_YELLOW;
                break;

            case STATE_YELLOW:
                gpio_pin_set_dt(&led0, 1);
                gpio_pin_set_dt(&led2, 1);
                current_state = STATE_RED;
                break;

            case STATE_RED:
                gpio_pin_set_dt(&led2, 1);
                current_state = STATE_GREEN;
                break;
        }

        k_msleep(SLEEP_TIME_MS);
    }

    return;
}
        \end{verbatim}
        \end{small}

    \section{Link do repositório}
        \url{https://github.com/KryptonPlusPlus/PSI3441/tree/main/atividade_01}

    
\end{document}
